% ===========================================================================
% OWNER: JOINT - lead M1 (Zoya Ahmad), written after chapters 1-10
% GRADED BY: Form 3 #1
% DRAW FROM: finished chapters only - no new claims here
% Rules: mid-eval-report/CLAUDE.md. Do not edit a file you do not own.
% ===========================================================================
\chapter{Conclusion}\label{ch:conclusion}
\paragraph{Problem.} When a Java test passes alone but fails after other tests, developers can see
that something is unstable but have no bounded, checkable way to find which earlier test caused it,
through which shared state, and how reliably (Section~\ref{ch1:sec:problem}). The three interviewed
practitioners do this by hand, and two of them said that an honest ``unknown'' is worth more than a
confident wrong answer.

\paragraph{Research findings.} Each part of the diagnosis has been solved on its own: delta
debugging reduces a failing order, RankF ranks likely polluters, and Takuan and PolDet explain
shared state. No system reviewed combines controlled reruns, a statistical reproduction rate,
code-level evidence and a typed refusal for one failure (Section~\ref{ch2:sec:gaps}). Shared state
alone is not proof: PolDet's own inspection judged 194 of its 324 reported polluters relevant
\cite{gyori2015poldet}. These findings fixed many of the design's rules and constants
(Table~\ref{ch2:tab:decisions}).

\paragraph{Requirements and design.} Seven functional requirements and the success criteria
(Section~\ref{ch1:sec:success}, Table~\ref{ch4:tab:fr}) define the product. A result is
\texttt{VERIFIED} only with a resource edge \emph{and} a Wilson lower bound of at least 0.70, which
takes at least 19 of 20 replays. The design is three modules joined by written contracts and one
report schema (Chapter~5). Four decisions shaped Iteration~1: static evidence with no
instrumentation (ADR-001), analysis depth that deepens only when needed (ADR-006), delta debugging
when no single polluter is found (ADR-007), and shuffled orders before declaring a failure not
reproduced (ADR-008).

\paragraph{Implementation and validation.} One command, \texttt{py -m runner diagnose}, now takes a
failing test through reproduction, the victim-alone check, polluter search, minimisation, repeated
verification, bytecode evidence and the outcome decision to a schema-validated report
(Chapter~6). On the seeded fixture, the four order-dependent cases are \texttt{VERIFIED}, each
reproducing 20 of 20 times and never alone, with the ground-truth polluters and resource; the three
negative controls end \texttt{UNRESOLVED} with a typed reason
(Table~\ref{ch8:tab:fixture-results}). 268 tests run in CI on every pull request
(Section~\ref{ch6:sec:m3}). These results come from hand-made fixture cases; on one real project, fastjson,
only the evidence extractor has been run (static analysis of its compiled classes; no fastjson
test has been run through the full tool), and user acceptance testing has not yet been run
(Section~\ref{ch8:sec:acceptance}).

\paragraph{Risks.} The main open risks are stated in the chapters they belong to: evidence from one
real project and a library that favours static analysis; no measured overhead for runtime evidence;
published baselines not yet run; no container yet, so diagnosis runs the project's tests as a local
process; and interviewed practitioners who do not use Java (Sections~\ref{ch2:sec:challenges}
and~\ref{ch1:sec:constraints}).

\paragraph{Continuation.} Iteration~2, before the FYP-1 Final evaluation, runs one real benchmark
project end to end, ranks candidates using the evidence module, makes a first container
deployment and adds runtime evidence where static evidence is not enough. Later iterations add the
execution budget, and FYP-2 adds filesystem evidence and the held-out evaluation against published
baselines (Section~\ref{ch3:sec:it2} and the Remaining FYP Plan).

Iteration~1 already shows the core of FlakeTrace's contribution on its fixture: for a known
order-dependent failure it returns either a result anyone can check (which tests, which shared
state, where in the code, how reliably) or a clear statement of why it cannot conclude.
